\documentclass[11pt]{amsart}
\usepackage{amsmath,amssymb,amsthm}
\usepackage[margin=1in]{geometry}
\usepackage{enumitem}
\usepackage{hyperref}

\newtheorem{lemma}{Lemma}
\newtheorem{theorem}{Theorem}
\renewcommand{\thetheorem}{\Alph{theorem}}
\theoremstyle{definition}
\newtheorem{definition}{Definition}
\theoremstyle{remark}
\newtheorem{remark}{Remark}

\newcommand{\Z}{\mathbb{Z}}
\newcommand{\Q}{\mathbb{Q}}
\newcommand{\C}{\mathbb{C}}
\newcommand{\HH}{\mathbb{H}}
\newcommand{\RP}{\mathbb{RP}}
\newcommand{\HP}{\mathbb{HP}}
\newcommand{\cA}{\mathcal{A}}
\newcommand{\cG}{\mathcal{G}}
\newcommand{\cL}{\mathcal{L}}
\newcommand{\B}{\mathfrak{B}}
\newcommand{\HPU}{\mathit{HPU}}
\newcommand{\cHPU}{\check{\mathit{HPU}}}
\DeclareMathOperator{\CS}{CS}
\DeclareMathOperator{\gr}{gr}
\DeclareMathOperator{\im}{im}
\DeclareMathOperator{\coker}{coker}
\DeclareMathOperator{\End}{End}
\DeclareMathOperator{\Map}{Map}

\title{A $PU(2)$-monopole Floer homology for integer homology spheres}
\date{}

\begin{document}

\begin{abstract}
We propose a definition with three ingredients: a real blow-up along
$\Phi=0$, on which the $SU(2)$-gauge action is free; Borel
$S^1$-equivariance with respect to scalar multiplication on $\Phi$; and
$\HPU(Y):=$ the $U$-torsion of the ``minus'' flavour, the Borel analogue
of $\im N$. Vanishing for $S^3$ is proved by a parity count, given the
standard analytic package. The instanton content is proved in the form of
a spectral sequence whose $E_1$-page, in high tower levels, is
$I_*(Y;\Q)$. Degeneration of that spectral sequence is not settled here.
\end{abstract}

\maketitle

\section{Setup}

Let $Y$ be an integer homology sphere with metric $g$, let $W\to Y$ be the
spinor bundle of the unique spin$^c$ structure with a fixed spin
connection, let $E=Y\times\C^2$, and let $\cG=\Map(Y,SU(2))$.
Configurations are pairs $(A,\Phi)$, where $A$ is an $SU(2)$-connection on
$E$ and $\Phi\in\Gamma(W\otimes E)$. The functional is
\[
\cL(A,\Phi)=\CS(A)+\tfrac12\langle D_A\Phi,\Phi\rangle_{L^2}.
\]
Its critical points satisfy
\[
\ast F_A=\rho^{-1}\big((\Phi\Phi^*)_{00}\big),\qquad D_A\Phi=0,
\]
with the sign chosen so that the Weitzenb\"ock term is
\[
\langle(\Phi\Phi^*)_{00}\Phi,\Phi\rangle=|(\Phi\Phi^*)_{00}|^2\ge\tfrac14|\Phi|^4.
\]
\begin{itemize}[leftmargin=*]
\item $S^1$ acts on $\Phi$ by scalars. This action is not a gauge symmetry.
\item Perturbations are $\cG$-invariant only; they need not be
  $S^1$-invariant.
\item Relative gradings are defined mod $6$: a degree-one gauge loop has
  index $8$ (ASD part) $-\,2$ (Dirac part, since $\int\mathrm{ch}_2=-1$)
  $=6$.
\end{itemize}

\begin{lemma}\label{lem:free}
$\cG$ acts freely on $\{\Phi\neq0\}$.
\end{lemma}

\begin{proof}
Suppose $g\cdot A=A$. Then $g$ is $A$-parallel, hence determined by its
value at one point. Suppose also $g\Phi=\Phi$, and pick $y$ with
$\Phi(y)\neq0$. Then $g(y)$ fixes the nonzero subspace
$\im\Phi(y)\subset E_y$. An element of $SU(2)$ with eigenvalue $1$ is the
identity, so $g(y)=1$, hence $g\equiv1$.
\end{proof}

Thus the only non-free points are those with $\Phi=0$. Irreducible flat
$\alpha$ have isotropy $\{\pm1\}$, acting on $\Phi$ by $-1$. The trivial
connection $\theta$ has isotropy $SU(2)$.

\section{Construction}

Form the real blow-up
\[
\B^\sigma=\{(A,r,\psi): r\ge0,\ \|\psi\|_{L^2}=1\}/\cG.
\]

\emph{Free quotient.} Lemma~\ref{lem:free} applies with $\psi$ in place of
$\Phi$, since $\psi\neq0$. Hence $\B^\sigma$ is a Hilbert manifold with
boundary $\partial\B^\sigma=\{r=0\}$.

\emph{KM flavours.} Kronheimer--Mrowka's construction applies to
$\B^\sigma$, giving $\check{\ }$, $\hat{\ }$, $\bar{\ }$ complexes with the
exact triangle. Since gradings are only $\Z/6$, one uses negative
completions.

\emph{Borel flavours.} The $S^1$-action on $\psi$ commutes with $\cG$.
Using $S^1$-equivariant families of perturbations parametrized by $ES^1$
yields $\check C^-$, $\check C^+$, $\check C^\infty$, related by the
triangle
\[
H^-\to H^\infty\to H^+\xrightarrow{\;N\;}H^-.
\]

\begin{definition}
\[
\HPU(Y):=\ker\big(\cHPU^-(Y)\to\cHPU^\infty(Y)\big)=\im N
=U\text{-torsion of }\cHPU^-(Y).
\]
\end{definition}

\section{Structure of the instanton face}

\begin{lemma}\label{lem:face}
\leavevmode
\begin{enumerate}[label=(\alph*)]
\item $\partial\B^\sigma=\big(\cA\times S(\Gamma(W\otimes E))\big)/\cG$,
  with $\cG$ acting freely and $S(\Gamma(W\otimes E))$ contractible. Hence
  $\partial\B^\sigma\simeq ESU(2)\times_{SU(2)}\widetilde{\B}_{\mathrm{inst}}$,
  the Borel model of the framed instanton configuration space.
\item The boundary critical points over irreducible $\alpha$ are
  $(\alpha,\pm\psi)$, with $\psi$ spanning a real eigenline of $L_\alpha$.
  They form $\RP^\infty$-towers with step $1$.
\item The boundary critical points over $\theta$ are $(\theta,E)$, with $E$
  a quaternionic eigenline of $L_\theta$. They form $\HP^\infty$-towers
  with step $4$; adjacent stable and unstable points differ in grading by
  $4-1=3$.
\item Every point $(\theta,E)$ is $S^1$-fixed for every perturbation,
  whereas $S^1/\{\pm1\}$ acts freely on the face over irreducible
  connections.
\end{enumerate}
\end{lemma}

\begin{proof}
(a) is immediate from Lemma~\ref{lem:free}. For (b),(c): $L_\theta$
commutes with $SU(2)$ acting on $\C^2=\HH$, so $L_\theta$ is
right-$\HH$-linear; its eigenspaces are quaternionic, generically of real
dimension $4$, and $SU(2)$ acts simply transitively on $S(E)=S^3$. For
(d): complex scalars commute with $SU(2)$, so lie in
$\End_{SU(2)}=\HH^{\mathrm{op}}$, and therefore preserve $E$. Over an
irreducible $\alpha$, $z\psi=\pm\psi$ only if $z=\pm1$.
\end{proof}

\section{Vanishing for $S^3$}

\begin{theorem}\label{thm:A}
$\HPU(S^3)=0$.
\end{theorem}

\begin{proof}
\emph{Step 1: the critical set is $\{\theta\}$.} Take the round metric.
For any critical point,
\[
0=\|\nabla\Phi\|^2+\int\tfrac{s}{4}|\Phi|^2+\|(\Phi\Phi^*)_{00}\|^2,
\]
so $\Phi=0$. Every flat $SU(2)$-connection on $S^3$ is trivial, so
$A=\theta$. It is nondegenerate: $H^1=0$, and $\ker D_\theta=0$ by
Lichnerowicz.

\emph{Step 2: small perturbations add nothing.} Take a small equivariant
family of perturbations. $\theta$ is critical for every $\cG$-invariant
perturbation, by $SU(2)$-isotropy. By compactness and the implicit
function theorem, the critical set of the family remains $\{\theta\}$.
Hence the generators of $\check C$ are exactly the boundary-stable points
$(\theta,E_j)$.

\emph{Step 3: every component of the differential vanishes.} Relative
gradings between boundary-stable $\theta$-points are $4m+6n$, which is
even. The equivariant components $\partial_k$ have degree $-1+2k$, which
is odd, so they vanish. The KM correction term
$\partial^u_s\bar\partial^s_u$ also vanishes: $\bar\partial^s_u$ requires
equal raw gradings between a stable and an unstable point, but these
differ by $4m+3+6n$, which is odd.

\emph{Step 4.} Hence $H^-=C^s\otimes\Q[[U]]$, which is $U$-torsion-free,
so $\HPU(S^3)=0$.
\end{proof}

\section{Instanton content over $\Q$}

\begin{theorem}\label{thm:B}
There is a spectral sequence whose $E_1$-page, in high tower levels, is
canonically $I_*(Y;\Q)$, regraded by the Dirac spectral flow $\sigma$, and
which converges to the high-level part of $\HPU(Y)$.
\end{theorem}

\begin{proof}
\emph{(i) Which contributions survive.} By Lemma~\ref{lem:face}(d), the
face over each irreducible $\alpha$ is $S^1$-free. Its Borel model is a
complex semi-infinite tower on which $U$ acts locally nilpotently, so it
is $U$-torsion and lands in $\HPU$. The $\theta$-points are $S^1$-fixed,
each contributing $\Q[[U]]$, which is torsion-free; so $\theta$ is
excluded. Free interior critical points (irreducible $PU(2)$ monopoles,
including the former Seiberg--Witten points, which are no longer special)
are $U$-torsion.

\emph{(ii) A filtration recovering $\partial_I$.} Work on the $\Z$-cover
and set $\nu=\gr_{PU}-\gr_I$.
\begin{itemize}[leftmargin=*]
\item A face component lying over an instanton trajectory of index $m$
  changes $\nu$ by $m-1\ge0$.
\item The $\nu$-preserving part is exactly Floer's differential
  $\partial_I$: each index-$1$ trajectory carries exactly one spinor line.
\item Components $\theta\to\beta$ raise $\nu$ by at least $3$. Hence
  $\theta$ spans a subcomplex of the associated graded, with quotient
  $CI\otimes(\text{tower})$.
\item At the bottoms of the half-towers, entries with negative Dirac
  spectral flow are truncated; for $\nu\gg0$ every entry is present.
\end{itemize}
Therefore, in high tower levels, $\gr_\nu$ of the $\alpha$-part is
canonically $(CI_*(Y;\Q),\partial_I)$, regraded by $\sigma$. This gives
the stated spectral sequence.
\end{proof}

\begin{remark}
The higher differentials come from Euler classes of the Dirac index
bundle over higher-index instanton moduli spaces, the analogue of the
$\mu$-class insertions in Feehan--Leness. Their vanishing is not proved
here.
\end{remark}

\section{Remarks}

\subsection*{Why these flavour choices}
\begin{itemize}[leftmargin=*]
\item The KM-reduced group $\coker i_*$ kills the $\alpha$-towers, just as
  it kills the reducible in KM theory, so it fails~(1).
\item The non-equivariant $\Q$-theory satisfies~(2) by the same parity
  argument, but fails~(1): the $\RP^\infty$-towers are $\Q$-acyclic, which
  deletes exactly the $\partial_I$-entries between $\alpha,\beta$ with
  nonzero Dirac spectral flow.
\item $S^1$-equivariance absorbs those spectral shifts into $U$-powers.
  The $U$-torsion of $H^-$ then separates the $S^1$-free part (the
  $\alpha$'s) from the $S^1$-fixed part ($\theta$).
\item The complement $U^{-1}H^-$ carries the $\theta$-towers. The torsion
  part also contains the free $PU(2)$ monopoles and the Seiberg--Witten
  data; this is where a Floer-theoretic Feehan--Leness relation would
  appear.
\end{itemize}

\subsection*{Analytic input assumed}
\begin{itemize}[leftmargin=*]
\item A $C^0$ bound on $\Phi$, from
  $|(\Phi\Phi^*)_{00}|^2\ge\tfrac14|\Phi|^4$.
\item Uhlenbeck strata of codimension $2\ell$: each bubble lowers the
  moduli dimension by $6$ and adds $4$ for its position, so these strata
  are irrelevant in the dimensions used.
\item Transversality for $ES^1$-families of perturbations.
\item KM-type gluing and completions.
\end{itemize}
Since the action on $\B^\sigma$ is free, no singular phenomena arise
beyond the boundary face already handled by Kronheimer--Mrowka.

\end{document}
